\documentclass[12pt]{article}
\usepackage[utf8]{inputenc}
\usepackage[T1]{fontenc}
\usepackage[margin=1in]{geometry}
\usepackage{setspace}
\usepackage{xcolor}
\usepackage[colorlinks=true, urlcolor=blue!60!black, linkcolor=blue!60!black]{hyperref}
\setlength{\parindent}{2em}
\setlength{\parskip}{1.25em}
\renewcommand{\baselinestretch}{1.0}
\begin{document}
\begin{center}
{\LARGE \textit{Students for Fair Admissions, Inc. v. President and Fellows of Harvard College}, 600 U.S. 181 (2023)\par}
\vspace{0.8em}
{\large Case Brief\par}
\end{center}
\vspace{1em}

\section{Memory Jogger}

The Supreme Court invalidated the race-conscious undergraduate admissions programs at Harvard College and the University of North Carolina under the Equal Protection Clause (and, via parallel standards, Title VI), holding that they lacked measurable objectives, used race as a negative and a stereotype, and had no logical endpoint---while allowing universities to consider an applicant's individual discussion of how race affected their life (\href{https://www.courtlistener.com/opinion/10049657/students-for-fair-admissions-inc-v-president-and-fellows-of-harvard/}{CourtListener}; \href{https://www.supremecourt.gov/opinions/22pdf/20-1199_hgdj.pdf}{Supreme Court Opinion}).

\section{Detailed Case Facts}

Petitioner Students for Fair Admissions, Inc. (SFFA) is a nonprofit membership organization founded in 2014 whose stated purpose is ``to defend human and civil rights secured by law, including the right of individuals to equal protection under the law'' (\href{https://www.law.cornell.edu/supremecourt/text/20-1199}{LII}). In November 2014, SFFA filed separate federal lawsuits challenging the race-conscious undergraduate admissions programs of two of the nation's oldest institutions of higher learning: Harvard College, a private university receiving federal funds and therefore subject to Title VI of the Civil Rights Act of 1964, 42 U.S.C. \S~2000d \textit{et seq.}, and the University of North Carolina at Chapel Hill (UNC), a public university subject directly to the Equal Protection Clause of the Fourteenth Amendment as well as Title VI (\href{https://www.law.cornell.edu/supremecourt/text/20-1199}{LII}; \href{https://www.supremecourt.gov/opinions/22pdf/600us1r53_4g15.pdf}{600 U.S. 181}).

\textbf{Harvard's Admissions Process:} Harvard evaluated tens of thousands of applicants through a multi-stage holistic review. First readers assigned numerical ratings across academic, extracurricular, athletic, school support, personal, and overall categories, and could take race into account when assigning the overall rating (\href{https://caselaw.findlaw.com/court/us-supreme-court/20-1199.html}{FindLaw}). Regional subcommittees and then the full 40-member admissions committee voted on applicants, periodically reviewing ``one-pagers'' displaying the racial breakdown of the tentatively admitted class relative to prior years (\href{https://caselaw.findlaw.com/court/us-supreme-court/20-1199.html}{FindLaw}; \href{https://directemployers.org/2023/07/03/scotus-ends-its-45-year-detour-around-the-equal-protection-clause-by-now-finding-race-based-admissions-practices-at-harvard-and-unc-unlawful/}{DirectEmployers}). In the final ``lop'' stage, when winnowing the tentative list to the target class size, the lop list displayed only four data points: legacy status, recruited athlete status, financial aid eligibility, and race (\href{https://caselaw.findlaw.com/court/us-supreme-court/20-1199.html}{FindLaw}). Across the admitted classes of 2009 to 2018, Black students remained within a narrow band of 10.0\% to 11.7\% of the admitted pool, and the District Court found that race served as a ``determinative tip'' for a significant share of admitted Black and Hispanic applicants while overall resulting in fewer Asian American and White students being admitted (\href{https://caselaw.findlaw.com/court/us-supreme-court/20-1199.html}{FindLaw}; \href{https://tikvah.org/wp-content/uploads/2026/05/Students-for-Fair-Admissions.pdf}{SFFA Opinion Excerpt}).

\textbf{UNC's Admissions Process:} At UNC, first readers evaluated applicants across more than forty criteria and were permitted to award a ``plus'' based on race, which the District Court found ``may be significant in an individual case'' (\href{https://caselaw.findlaw.com/court/us-supreme-court/20-1199.html}{FindLaw}). First-reader decisions were provisionally final and then reviewed by a ``school group review'' committee that could also consider the applicant's race (\href{https://caselaw.findlaw.com/court/us-supreme-court/20-1199.html}{FindLaw}). Both Harvard and UNC defended their consideration of race as necessary to achieve the educational benefits of student-body diversity previously recognized in \textit{Regents of the University of California v. Bakke}, 438 U.S. 265 (1978), \textit{Grutter v. Bollinger}, 539 U.S. 306 (2003), and \textit{Fisher v. University of Texas at Austin}, 579 U.S. 365 (2016) (\href{https://caselaw.findlaw.com/court/us-supreme-court/20-1199.html}{FindLaw}).

\section{Procedural History}

\begin{enumerate}
\item \textbf{U.S. District Court for the District of Massachusetts (2019):} Following a 15-day bench trial featuring testimony from 30 witnesses, Judge Allison D. Burroughs upheld Harvard's admissions program under Title VI and existing Supreme Court precedent, finding no discriminatory animus or unlawful quotas. \textit{Students for Fair Admissions, Inc. v. President \& Fellows of Harvard College}, 397 F. Supp. 3d 126 (D. Mass. 2019) (\href{https://www.law.cornell.edu/supremecourt/text/20-1199}{LII}).
\item \textbf{U.S. Court of Appeals for the First Circuit (2020):} The First Circuit affirmed the District Court's judgment in full, holding that Harvard's limited use of race survived strict scrutiny under \textit{Grutter} and \textit{Fisher}. 980 F.3d 157 (1st Cir. 2020) (\href{https://www.law.cornell.edu/supremecourt/text/20-1199}{LII}).
\item \textbf{U.S. District Court for the Middle District of North Carolina (2021):} After an eight-day bench trial, the District Court upheld UNC's admissions program under the Equal Protection Clause and Title VI, concluding that race-neutral alternatives were not yet workable. 567 F. Supp. 3d 580 (M.D.N.C. 2021) (\href{https://www.law.cornell.edu/supremecourt/text/20-1199}{LII}; \href{https://www.ncmd.uscourts.gov/students-fair-admissions-inc-v-president-fellows-harvard-college}{M.D.N.C.}).
\item \textbf{U.S. Supreme Court (Argued October 31, 2022; Decided June 29, 2023):} The Court granted certiorari in \textit{Harvard} (No. 20-1199) and certiorari before judgment in \textit{UNC} (No. 21-707), later unconsolidating them for oral argument so Justice Ketanji Brown Jackson could recuse herself from \textit{Harvard} (having served on the Harvard Board of Overseers) while participating in \textit{UNC} (\href{https://www.scotusblog.com/cases/students-for-fair-admissions-inc-v-president-fellows-of-harvard-college/}{SCOTUSblog (Harvard)}; \href{https://www.scotusblog.com/cases/students-for-fair-admissions-inc-v-university-of-north-carolina/}{SCOTUSblog (UNC)}). On June 29, 2023, the Court issued a single consolidated opinion reversing both lower courts (\href{https://www.supremecourt.gov/opinions/22pdf/20-1199_hgdj.pdf}{Supreme Court Opinion}).
\end{enumerate}

\section{Judicial Votes}

\textbf{Vote Breakdown:} \textbf{6--2} in \textit{Harvard} (No. 20-1199) and \textbf{6--3} in \textit{UNC} (No. 21-707) (\href{https://www.law.cornell.edu/supremecourt/text/20-1199}{LII}; \href{https://www.oyez.org/cases/2022/20-1199}{Oyez}).
\begin{itemize}
\item \textbf{Majority (6):} Chief Justice John Roberts (author), joined by Justices Clarence Thomas, Samuel Alito, Neil Gorsuch, Brett Kavanaugh, and Amy Coney Barrett (\href{https://www.law.cornell.edu/supremecourt/text/20-1199}{LII}).
\item \textbf{Concurrences:} Justice Clarence Thomas filed a concurring opinion; Justice Neil Gorsuch filed a concurring opinion, in which Justice Thomas joined; and Justice Brett Kavanaugh filed a concurring opinion (\href{https://www.law.cornell.edu/supremecourt/text/20-1199}{LII}).
\item \textbf{Dissents:} Justice Sonia Sotomayor filed a dissenting opinion, in which Justice Elena Kagan joined in full and Justice Ketanji Brown Jackson joined as it applies to \textit{UNC} (No. 21-707). Justice Ketanji Brown Jackson filed a separate dissenting opinion in \textit{UNC} (No. 21-707), in which Justices Sotomayor and Kagan joined (\href{https://www.law.cornell.edu/supremecourt/text/20-1199}{LII}).
\item \textbf{Recusal:} Justice Ketanji Brown Jackson took no part in the consideration or decision of \textit{Harvard} (No. 20-1199) (\href{https://www.law.cornell.edu/supremecourt/text/20-1199}{LII}).
\end{itemize}

\section{Holding}

\textbf{Judgment:} Reversed in both No. 20-1199 (\textit{Harvard}) and No. 21-707 (\textit{UNC}) (\href{https://www.law.cornell.edu/supremecourt/text/20-1199}{LII}).

\textbf{Rule Announced:} First, SFFA satisfied Article III organizational standing under \textit{Hunt v. Washington State Apple Advertising Commission}, 432 U.S. 333 (1977) (\href{https://caselaw.findlaw.com/court/us-supreme-court/20-1199.html}{FindLaw}). Second, on the merits, Harvard's and UNC's race-based admissions programs violate the Equal Protection Clause of the Fourteenth Amendment (and, because the parties agreed Title VI tracks the Equal Protection Clause, Title VI as well) because they fail strict scrutiny: their asserted diversity objectives are too amorphous to permit meaningful judicial review, they unavoidably employ race as a ``negative'' in a zero-sum admissions process, they rely on impermissible racial stereotyping, and they lack any logical endpoint (\href{https://www.law.cornell.edu/supremecourt/text/20-1199}{LII}). At the same time, universities remain free to consider an applicant's discussion of how race affected their own life, so long as any resulting benefit is concretely tied to a quality of character or unique ability the individual can contribute (\href{https://www.law.cornell.edu/supremecourt/text/20-1199}{LII}).

\textbf{Holding vs. Dicta:}
\begin{itemize}
\item \textbf{Strict Holdings:} (1) SFFA had Article III associational standing; and (2) Harvard's and UNC's race-conscious admissions systems are unconstitutional under the Equal Protection Clause (and unlawful under Title VI) because they fail strict scrutiny's requirements of measurable compelling interests, narrow tailoring without negative racial impacts or stereotyping, and a defined termination point (\href{https://www.law.cornell.edu/supremecourt/text/20-1199}{LII}).
\item \textbf{Important Carve-Outs and Dicta:}
  \begin{itemize}
  \item \textit{Military Academies (Footnote 4):} Because no U.S. military academy was a party and lower courts had not addressed them, the majority explicitly reserved whether race-based admissions at service academies might be justified by ``potentially distinct interests'' (\href{https://www.lawfaremedia.org/article/the-national-security-footnote-in-the-affirmative-action-decision}{Lawfare}).
  \item \textit{Personal-Experience Essays:} The Court's statement permitting consideration of an applicant's essay discussing how race affected their life---while cautioning that universities may not ``simply establish through application essays or other means the regime we hold unlawful today''---guides future litigation without resolving specific essay prompts (\href{https://www.forbes.com/sites/michaeltnietzel/2023/07/02/three-of-the-bigger-puzzles-in-the-supreme-court-college-admissions-decision/}{Forbes}; \href{https://www.law.cornell.edu/supremecourt/text/20-1199}{LII}).
  \item \textit{Status of Grutter:} Although the majority did not formally state that it was overruling \textit{Grutter v. Bollinger}, Justices Thomas and Sotomayor both observed that the decision effectively ended race-based affirmative action in civilian higher education (\href{https://www.law.cornell.edu/supremecourt/text/20-1199}{LII}; \href{https://en.wikisource.org/wiki/Students_for_Fair_Admissions_v._President_and_Fellows_of_Harvard_College/Opinion_of_Justice_Thomas}{Wikisource}).
  \end{itemize}
\end{itemize}

\section{Analysis of Opinions}

\subsection{Majority Opinion (Chief Justice Roberts)}

Chief Justice Roberts began by confirming SFFA's organizational standing under \textit{Hunt}, rejecting UNC's contention that SFFA was not a ``genuine'' membership organization (\href{https://www.law.cornell.edu/supremecourt/text/20-1199}{LII}). On the merits, Roberts traced Equal Protection doctrine from the ratification of the Fourteenth Amendment through \textit{Strauder v. West Virginia}, \textit{Yick Wo v. Hopkins}, and \textit{Brown v. Board of Education}, concluding that the Clause's core purpose is universal and colorblind: ``Eliminating racial discrimination means eliminating all of it'' (\href{https://caselaw.findlaw.com/court/us-supreme-court/20-1199.html}{FindLaw}; \href{https://directemployers.org/2023/07/03/scotus-ends-its-45-year-detour-around-the-equal-protection-clause-by-now-finding-race-based-admissions-practices-at-harvard-and-unc-unlawful/}{DirectEmployers}). Any racial classification must survive strict scrutiny---advancing a compelling governmental interest through narrowly tailored (necessary) means---and outside of remedying specific, identified instances of past discrimination or avoiding imminent prison violence, the Court has almost never tolerated race-based state action (\href{https://caselaw.findlaw.com/court/us-supreme-court/20-1199.html}{FindLaw}).

Measuring Harvard's and UNC's programs against \textit{Grutter}'s guardrails, Roberts identified four fatal defects. First, the universities' stated educational goals (such as training future leaders, promoting robust exchange of ideas, and enhancing appreciation of differences) were ``commendable'' but too elusive and unmeasurable for courts to verify when they are achieved, and the six broad racial categories employed were themselves imprecise, overbroad, and underinclusive (\href{https://www.ncmd.uscourts.gov/students-fair-admissions-inc-v-president-fellows-harvard-college}{M.D.N.C.}; \href{https://tikvah.org/wp-content/uploads/2026/05/Students-for-Fair-Admissions.pdf}{SFFA Opinion Excerpt}). Second, because ``[c]ollege admissions are zero-sum,'' awarding a racial ``plus'' to some applicants necessarily operated as a ``negative'' against others, including Asian American applicants (\href{https://tikvah.org/wp-content/uploads/2026/05/Students-for-Fair-Admissions.pdf}{SFFA Opinion Excerpt}). Third, assuming that an applicant's race predicts a distinct viewpoint or contribution engaged in forbidden racial stereotyping (\href{https://caselaw.findlaw.com/court/us-supreme-court/20-1199.html}{FindLaw}). Fourth, neither university offered a meaningful endpoint, instead maintaining stable annual racial percentages ``like clockwork'' and relying on periodic reviews that \textit{Grutter} never treated as a substitute for a sunset (\href{https://www.law.cornell.edu/supremecourt/text/20-1199}{LII}; \href{https://directemployers.org/2023/07/03/scotus-ends-its-45-year-detour-around-the-equal-protection-clause-by-now-finding-race-based-admissions-practices-at-harvard-and-unc-unlawful/}{DirectEmployers}).

\subsection{Concurring Opinion (Justice Thomas)}

Justice Thomas wrote separately to offer an originalist defense of a strictly ``colorblind Constitution'' rooted in the Declaration of Independence, the Civil Rights Act of 1866, and the Fourteenth Amendment's original public meaning (\href{https://en.wikisource.org/wiki/Students_for_Fair_Admissions_v._President_and_Fellows_of_Harvard_College/Opinion_of_Justice_Thomas}{Wikisource}). Directly answering the dissents' reliance on Reconstruction-era legislation such as the Freedmen's Bureau Acts, Thomas argued that those statutes were facially race-neutral measures aiding formerly enslaved people and destitute refugees rather than race-based preferences (\href{https://en.wikisource.org/wiki/Students_for_Fair_Admissions_v._President_and_Fellows_of_Harvard_College/Opinion_of_Justice_Thomas}{Wikisource}). He also rejected deference to universities' self-assessments of diversity, argued that \textit{Grutter} was ``for all intents and purposes, overruled,'' and warned that race-conscious admissions stigmatize beneficiaries and undermine individual dignity (\href{https://en.wikisource.org/wiki/Students_for_Fair_Admissions_v._President_and_Fellows_of_Harvard_College/Opinion_of_Justice_Thomas}{Wikisource}).

\subsection{Concurring Opinion (Justice Gorsuch, Joined by Justice Thomas)}

Justice Gorsuch focused on the plain text of Title VI of the Civil Rights Act of 1964, 42 U.S.C. \S~2000d, which prohibits any federally funded program from discriminating against any person ``on the ground of race, color, or national origin'' (\href{https://www.law.cornell.edu/supremecourt/text/20-1199}{LII}; \href{https://www.supremecourt.gov/opinions/22pdf/20-1199_hgdj.pdf}{Supreme Court Opinion}). Drawing an analogy to his textualist analysis of Title VII in \textit{Bostock v. Clayton County}, Gorsuch argued that treating an individual worse than others because of race violates Title VI's categorical text regardless of benign motives or constitutional strict scrutiny, and he criticized the federal government's six racial check-box categories as arbitrary and incoherent (\href{https://www.law.cornell.edu/supremecourt/text/20-1199}{LII}; \href{https://tikvah.org/wp-content/uploads/2026/05/Students-for-Fair-Admissions.pdf}{SFFA Opinion Excerpt}).

\subsection{Concurring Opinion (Justice Kavanaugh)}

Justice Kavanaugh wrote briefly to emphasize that the Court's ruling followed directly from \textit{Grutter v. Bollinger}'s own 25-year temporal limit, under which Justice O'Connor had expected in 2003 that race-conscious admissions would no longer be necessary by 2028 (\href{https://www.supremecourt.gov/opinions/22pdf/20-1199_hgdj.pdf}{Supreme Court Opinion}). Because Harvard's and UNC's programs had no sunset and would extend racial preferences indefinitely beyond the 2028 horizon, Kavanaugh concluded they could not be reconciled with precedent, while noting that race-neutral socioeconomic and geographic outreach remains available (\href{https://www.law.cornell.edu/supremecourt/text/20-1199}{LII}).

\subsection{Dissenting Opinion (Justice Sotomayor, Joined by Justice Kagan, and by Justice Jackson as to No. 21-707)}

Justice Sotomayor accused the majority of ``cement[ing] a superficial rule of colorblindness as a constitutional principle in an endemically segregated society'' and rolling back 45 years of settled precedent from \textit{Bakke} through \textit{Fisher} (\href{https://www.quinnemanuel.com/the-firm/publications/the-road-ahead-fostering-diversity-after-the-supreme-court-s-decision-overturning-affirmative-action-in-university-admissions/}{Quinn Emanuel}; \href{https://directemployers.org/2023/07/03/scotus-ends-its-45-year-detour-around-the-equal-protection-clause-by-now-finding-race-based-admissions-practices-at-harvard-and-unc-unlawful/}{DirectEmployers}). Invoking the Reconstruction Congress's enactment of the Freedmen's Bureau Acts of 1865 and 1866 alongside the Fourteenth Amendment, she argued that the Equal Protection Clause was never understood to forbid race-conscious remedies aimed at dismantling racial caste and achieving genuine equal opportunity (\href{https://www.law.cornell.edu/supremecourt/text/20-1199}{LII}). Sotomayor also defended the District Courts' extensive factual findings that neither Harvard nor UNC used quotas or discriminated against Asian American applicants, and she criticized the majority's footnote 4 carve-out for military academies as revealing the inconsistency of denying civilian universities the same diversity interest (\href{https://www.lawfaremedia.org/article/the-national-security-footnote-in-the-affirmative-action-decision}{Lawfare}; \href{https://www.law.cornell.edu/supremecourt/text/20-1199}{LII}).

\subsection{Dissenting Opinion (Justice Jackson, Joined by Justices Sotomayor and Kagan, in No. 21-707)}

Writing solely in \textit{UNC}, Justice Jackson grounded her dissent in the historical and empirical reality of ``gulf-sized race-based gaps'' in American wealth, health, housing, and education---citing, for example, 2019 data showing median White family wealth of roughly \$188,000 compared to \$24,000 for Black families (\href{https://www.quinnemanuel.com/the-firm/publications/the-road-ahead-fostering-diversity-after-the-supreme-court-s-decision-overturning-affirmative-action-in-university-admissions/}{Quinn Emanuel}). Tracing intergenerational disparities from slavery and Jim Crow through federal housing and GI Bill exclusion, Jackson argued that UNC's holistic review did not perpetuate racial hierarchy but rather accounted for the real-world obstacles applicants faced. She sharply criticized the majority's ``let-them-eat-cake obliviousness'' in declaring race irrelevant by legal fiat while leaving legacy and donor preferences intact (\href{https://www.quinnemanuel.com/the-firm/publications/the-road-ahead-fostering-diversity-after-the-supreme-court-s-decision-overturning-affirmative-action-in-university-admissions/}{Quinn Emanuel}).

\section{Examples: Future Applications}

\begin{enumerate}
\item \textbf{Same-Side (State Medical School Tracking Racial Target Bands in Holistic Review):} A public medical school uses a holistic admissions rubric in which committee members receive weekly demographic dashboards and apply an informal ``plus'' factor to maintain Black and Hispanic enrollment within a steady 12\%--14\% band each year without any sunset date. Under \textit{SFFA}, this program is plainly unconstitutional: it replicates the exact defects condemned in \textit{Harvard} and \textit{UNC}---unmeasurable diversity goals, zero-sum racial preferences, stereotyping, and no logical endpoint (\href{https://www.law.cornell.edu/supremecourt/text/20-1199}{LII}).

\item \textbf{Same-Side (Private University Giving Automatic Points for Checking a Racial Box on Scholarship Applications):} A federally funded private university awards automatic merit-scholarship points to applicants who check specific underrepresented racial boxes to achieve ``campus demographic parity.'' Because Title VI tracks the Equal Protection Clause under \textit{SFFA} (and under Justice Gorsuch's textualist reading of 42 U.S.C. \S~2000d), awarding competitive financial benefits based solely on racial classification is unlawful (\href{https://www.law.cornell.edu/supremecourt/text/20-1199}{LII}).

\item \textbf{Opposite-Side (Individualized Essay Credit for Overcoming Specific Racial Discrimination):} An applicant to a public law school writes a personal statement describing how she organized a student legal-aid clinic after confronting racial discrimination in her hometown's housing market, and the admissions committee awards her credit for demonstrated leadership, resilience, and initiative. Under Chief Justice Roberts's express safe harbor in \textit{SFFA}, this is constitutional because the benefit is tied concretely to the applicant's individual character and unique ability to contribute rather than to race \textit{per se} (\href{https://www.law.cornell.edu/supremecourt/text/20-1199}{LII}; \href{https://www.forbes.com/sites/michaeltnietzel/2023/07/02/three-of-the-bigger-puzzles-in-the-supreme-court-college-admissions-decision/}{Forbes}).

\item \textbf{Opposite-Side (Race-Neutral Socioeconomic and First-Generation Admissions Preferences):} A selective state university eliminates all access to applicants' racial check-boxes and instead grants an admissions preference to low-income, first-generation college applicants from under-resourced ZIP codes, applied identically to students of every race. Because the classification is facially race-neutral and not adopted with racially discriminatory animus, it is permissible under \textit{SFFA} and does not trigger strict scrutiny (\href{https://www.lawyerscommittee.org/wp-content/uploads/2023/10/20220927-LC_QA_Affirmative-Action_FINAL.pdf}{Lawyers' Committee for Civil Rights}).

\item \textbf{Fence-Sitter (Race-Conscious Admissions at U.S. Military Service Academies):} The U.S. Military Academy at West Point considers race in awarding ``letters of assurance'' and assembling its incoming officer class, asserting a national-security interest in maintaining a racially diverse officer corps to ensure unit cohesion and legitimacy. This is a quintessential fence-sitter: \textit{SFFA}'s Footnote 4 expressly left military academies undecided due to ``potentially distinct interests,'' sparking immediate follow-on litigation (\href{https://law.justia.com/cases/federal/district-courts/new-york/nysdce/7:2023cv08262/606600/78/}{S.D.N.Y. 2024}; \href{https://www.lawfaremedia.org/article/the-national-security-footnote-in-the-affirmative-action-decision}{Lawfare}), yet critics note that most military officers are commissioned through civilian university ROTC programs that are now barred from using race (\href{https://www.supremecourt.gov/DocketPDF/23/23A696/298928/20240126173859280_Army\%20SCOTUS\%20App.pdf}{West Point Emergency Application}).
\end{enumerate}

\section{Critique}

\textbf{Originalist and Textualist Perspectives:} Conservative originalists and textualists broadly applaud \textit{SFFA} for restoring the anti-classification principle of the Equal Protection Clause and the plain text of Title VI, arguing that \textit{Grutter} had created an unworkable doctrinal exception that invited judicial deference to racial balancing and entrenched arbitrary bureaucratic racial categories (\href{https://directemployers.org/2023/07/03/scotus-ends-its-45-year-detour-around-the-equal-protection-clause-by-now-finding-race-based-admissions-practices-at-harvard-and-unc-unlawful/}{DirectEmployers}; \href{https://dldlawyers.com/justice-clarence-thomas-on-affirmative-action-and-equal-protection-under-the-law/}{DLD Analysis}). At the same time, even some originalist commentators note internal tensions in the majority's work: Chief Justice Roberts declined to rely primarily on Fourteenth Amendment original public meaning (leaving that debate to Justice Thomas), avoided formally overruling \textit{Grutter}, and carved out military academies in Footnote 4 without explaining why strict colorblindness should yield to national-security deference (\href{https://www.lawfaremedia.org/article/the-national-security-footnote-in-the-affirmative-action-decision}{Lawfare}; \href{https://mindingthecampus.org/2023/07/14/sffa-v-harvard-should-our-nations-service-academies-be-exempt-from-this-landmark-ruling/}{Minding the Campus}).

\textbf{Progressive and Doctrinal Perspectives:} Progressive legal scholars and historians sharply challenge both the majority's history and its doctrinal consequences. Critics argue that the Reconstruction Congress's enactment of race-conscious Freedmen's Bureau legislation refutes Justice Thomas's claim that the Fourteenth Amendment originally mandated strict colorblindness, contending instead that the Amendment was designed around an \textit{anti-subordination} principle to uplift Black Americans and dismantle racial caste (\href{https://aulawreview.org/wp-content/uploads/2023/09/8_HinojosaTorres_Print.pdf}{American University Law Review}; \href{https://www.albanylawreview.org/article/116873-supreme-mirage-clarence-thomas-s-invented-history-of-colorblindness-originalism/attachment/250380.pdf}{Albany Law Review}). Doctrinally, critics fault the Court for disregarding detailed district-court factual findings, destabilizing 45 years of reliance interests across higher education, threatening declines in Black and Hispanic representation in selective universities and the legal profession, and spurring aggressive challenges to corporate diversity programs and race-neutral admissions reforms (\href{https://digitalcommons.law.seattleu.edu/cgi/viewcontent.cgi?article=2876&context=sulr}{Seattle University Law Review}; \href{https://verdict.justia.com/2023/10/11/future-of-affirmative-action-in-employment-decisions}{Justia's Verdict}; \href{https://law.justia.com/cases/federal/appellate-courts/ca5/24-50631/24-50631-2025-07-11.html}{5th Cir. 2025}).

\textbf{Underlying Competing Values:} \textit{SFFA} crystallizes the deepest philosophical divide in American Equal Protection law: \textit{anti-classification (formal colorblindness)}---which insists that the Constitution protects individuals rather than groups and forbids the state from sorting citizens by race regardless of benevolent purpose---versus \textit{anti-subordination (substantive equality)}---which insists that the Fourteenth Amendment permits race-conscious remedies to overcome enduring structural inequities rooted in centuries of state-sponsored racial exclusion.

\section{Key Quotations}

\begin{quote}
``Eliminating racial discrimination means eliminating all of it.''
\end{quote}
\noindent --- Chief Justice John Roberts, Opinion of the Court, 600 U.S. at 206 (\href{https://www.supremecourt.gov/opinions/22pdf/20-1199_hgdj.pdf}{Supreme Court Opinion}; \href{https://www.courtlistener.com/opinion/10049657/students-for-fair-admissions-inc-v-president-and-fellows-of-harvard/}{CourtListener}).

\begin{quote}
``College admissions are zero-sum. A benefit provided to some applicants but not to others necessarily advantages the former group at the expense of the latter.''
\end{quote}
\noindent --- Chief Justice John Roberts, Opinion of the Court, 600 U.S. at 218--219 (\href{https://tikvah.org/wp-content/uploads/2026/05/Students-for-Fair-Admissions.pdf}{SFFA Opinion Excerpt}; \href{https://www.law.cornell.edu/supremecourt/text/20-1199}{LII}).

\begin{quote}
``At the same time, nothing prohibits universities from considering an applicant's discussion of how race affected the applicant's life, so long as that discussion is concretely tied to a quality of character or unique ability that the particular applicant can contribute to the university.''
\end{quote}
\noindent --- Chief Justice John Roberts, Opinion of the Court, 600 U.S. at 230--231 (\href{https://www.law.cornell.edu/supremecourt/text/20-1199}{LII}).

\begin{quote}
``[T]he Court cements a superficial rule of colorblindness as a constitutional principle in an endemically segregated society . . . . In a society where race continues to matter, there is no constitutional requirement that institutions attempting to remedy their legacies of racial exclusion must operate with a blindfold.''
\end{quote}
\noindent --- Justice Sonia Sotomayor, Dissenting Opinion, 600 U.S. at 318, 383 (\href{https://www.quinnemanuel.com/the-firm/publications/the-road-ahead-fostering-diversity-after-the-supreme-court-s-decision-overturning-affirmative-action-in-university-admissions/}{Quinn Emanuel}; \href{https://www.law.cornell.edu/supremecourt/text/20-1199}{LII}).

\begin{quote}
``With let-them-eat-cake obliviousness, today, the majority pulls the ripcord and announces `colorblindness for all' by legal fiat. But deeming race irrelevant in law does not make it so in life.''
\end{quote}
\noindent --- Justice Ketanji Brown Jackson, Dissenting Opinion (No. 21-707), 600 U.S. at 407 (\href{https://www.quinnemanuel.com/the-firm/publications/the-road-ahead-fostering-diversity-after-the-supreme-court-s-decision-overturning-affirmative-action-in-university-admissions/}{Quinn Emanuel}; \href{https://www.law.cornell.edu/supremecourt/text/20-1199}{LII}).

\end{document}
